\newcommand{\feat}{\texttt}

%\paragraph*{Abstract}
%\begin{abstract}
%	TODO
%\end{abstract}



%      \begin{figure}[!htb]
%        \center{\includegraphics[width=0.5\textwidth]
%        {images/20190419001057_1.jpg}}
%        \caption{\label{fig1} Player performance, accessable in-game: Number of normal attacks, number of CA, number of CA-Chips, etc.}
%      \end{figure}

%  (fig. \ref{fig1}, \ref{fig2}) 
%  An example is Chitnis et al.\cite{chitnis} 


\section{Motivation}
Object detection or object recognition is one of the most common problem across scientific fields. One of the more popular In the scope of this project we want to implement a 


\section{Data}
The dataset consists of 17,500 aerial images out of which 13,136 contain an columnar cacti  and 4,364 do not.  Each image is 32x32 in size. The images have been resized by kaggle to be uniform in size. The dataset also contains an .csv file that annoates for each image if it contains an columnar cacti or not. The images are all from the Tehucan-Cuicatlan valley in the south of Mexico.  \cite{LOPEZJIMENZ2019} The images have been obtained using a drone from an flight altitude of 100 m. The images were manually labeled.
      \begin{figure}[!htb]
        \center{\includegraphics[width=0.5\textwidth]
        {images/test.jpg}}
        \caption{\label{fig1} test \cite{whit}}
      \end{figure}